% -------------
\section{Domain-Aware Parameterization}
\label{sec:appendix_param_details}

This section details the sampling constraints of~\autoref{subsec:templates}: templates constrain
the values they sample in four ways, shown below with examples from all five branches.
\autoref{tab:authoritative_sources} lists the sources of the tabulated values.

\paragraph{Tabulated Properties.} A template selects a substance, material, section or scenario
from a curated table and takes the properties the problem needs from that entry, so that they are
mutually consistent.
Examples are the critical properties, heats of formation and heat-capacity coefficients
($C_p/R = A + BT + CT^2 + DT^{-2}$) of chemical species; the Young's modulus, Poisson's ratio and
shear modulus of materials such as 6061-T6 aluminum, in SI or US customary units; the properties
that set the phase velocity in media such as fused silica; the depth, area and moment of inertia of
AISC wide-flange sections; the ranges of permeability and friction angle for each soil type; and the
ranges of arrival and service rates for each queueing scenario.

\paragraph{Laws, Correlations and Standards.} Other values follow the relations and design tables
used in practice.
Reaction problems use balanced chemical equations, compression indices follow Skempton's
correlation with the liquid limit, and bearing-capacity factors are Terzaghi's values for the
sampled friction angle.
Manning's coefficients follow the channel lining, runoff coefficients the land use, and SCS curve
numbers the land use and hydrologic soil group.
Acceptance-sampling plans come from the tables of MIL-STD-105E, control-chart factors from their
statistical definitions, and safety factors from standard normal quantiles at typical
cycle-service levels.

\paragraph{Regime-Aware Ranges.} Where the regime matters, templates set sampling ranges relative to
the quantity that decides it, so that each instance stays in the regime its problem assumes or lands
in the one it tests.
A downsampling problem draws the signal frequency relative to the aliasing limit, so that the result
aliases or not, and a forced-vibration problem draws the forcing frequency relative to the natural
frequency, so that instances fall below, near and above resonance.
Channel slopes keep the flow in the regime the problem assumes, such as subcritical uniform flow;
beam loads keep deflections within serviceability limits without overstressing the section; queue
utilizations stay below one, so that a steady state exists; and production rates are 1.5 to 6 times
the demand rate, so that finite-rate production is feasible.
Some templates also check each drawn scenario and redraw it when it leaves the regime: the
falling-film template, for example, redraws until the film Reynolds number is below 1,500 and the
maximum velocity below 10\,m/s.

\paragraph{Exact Values.} Templates draw some values from structured sets, so that they have the
form the problem requires and the arithmetic stays exact.
Probabilities of discrete random variables sum to exactly one as printed, rows of Markov transition
matrices sum exactly to one, modulation orders are powers of two, carrier frequencies are integer
multiples of the bit rate, and discrete-time frequencies are rational multiples of $\pi$.


% -------------
\section{Template Implementation Examples}
\label{sec:appendix_template_examples}

The three listings below show templates as they stand in the code, one per level and from three
branches, without their documentation strings.
The Easy template (natural frequency of a torsional pendulum, mechanical engineering) gives the
disc's mass moment of inertia and the shaft's torsional stiffness, sampled over ranges from small
rotors to flywheels, and asks for the natural frequency and period without stating the formula that
relates them.
The Intermediate template (mean time to failure of a cooling loop, industrial engineering) samples
the system topology: for pumps in series, the lifetime of the system is the minimum of independent
exponential lifetimes, whereas for two pumps in parallel, the memoryless property splits the
lifetime into two phases, so a sampled parameter changes the derivation, not only the numbers.
The Advanced template (molar volumes from the van der Waals equation, chemical engineering) couples
the equation's parameters, computed from the critical properties, with the cubic equation of state,
which it solves numerically for three real roots at a temperature and pressure that it redraws until
the state lies in the two-phase region; the gold trace of each of its instances has six or more
milestones.
Every template prints each value at a stated precision, and every printed calculation recomputes
from its printed operands (\autoref{appendix:certification}); the templates whose procedures round
at a stated precision in practice also prescribe the rounding of their answer
(\autoref{appendix:scoring}).

\vspace{6pt}

\subsection{Easy}

\vspace{6pt}

\begin{lstlisting}[style=pythonstyle]
def template_undamped_natural_frequency_torsional():
    # 1. Parameterize the inputs with random values

    # Mass moment of inertia (J_0): Represents the rotational inertia of the disc.
    # The range covers small rotors to medium-sized flywheels.
    mass_moment_of_inertia = round(random.uniform(0.05, 25.0), 3)  # in kg-m^2

    # Torsional stiffness (k_t): Represents the shaft's resistance to twisting.
    torsional_stiffness = random.randint(200, 75000)                # in N-m/rad

    # Standardize precision for all calculations and final outputs
    precision = 3

    # 2. Perform the core calculations for the solution

    # Step A: Calculate the undamped natural frequency in radians per second
    omega_n = math.sqrt(torsional_stiffness / mass_moment_of_inertia)

    # Step B: Convert the natural frequency from rad/s to Hertz (Hz)
    f_n = omega_n / (2 * math.pi)

    # Step C: Calculate the natural period of oscillation
    tau_n = 1 / f_n

    # 3. Generate the question and solution strings

    question = (
        f"A torsional pendulum consists of a disc with a mass moment of inertia of "
        f"{mass_moment_of_inertia} kg-m^2 attached to a shaft with a torsional stiffness of "
        f"{torsional_stiffness} N-m/rad. The system is undamped. "
        f"Calculate the natural frequency of torsional vibration in both rad/s and Hz, "
        f"and determine the corresponding period."
    )

    solution = (
        f"**Given:**\n"
        f"Mass Moment of Inertia (J_0) = {mass_moment_of_inertia} kg-m^2\n"
        f"Torsional Stiffness (k_t) = {torsional_stiffness} N-m/rad\n\n"

        f"**Step 1:** Calculate the undamped natural frequency in radians per second (omega_n).\n"
        f"The formula for a torsional system is: omega_n = sqrt(k_t / J_0)\n"
        f"omega_n = sqrt({torsional_stiffness} / {mass_moment_of_inertia})\n"
        f"omega_n = {round(omega_n, precision)} rad/s\n\n"

        f"**Step 2:** Convert the natural frequency to Hertz (f_n).\n"
        f"The relationship is: f_n = omega_n / (2 * pi)\n"
        f"f_n = {round(omega_n, precision)} / (2 * pi)\n"
        f"f_n = {round(f_n, precision)} Hz\n\n"

        f"**Step 3:** Calculate the natural period of vibration (tau_n).\n"
        f"The period is the reciprocal of the frequency in Hz.\n"
        f"The formula is: tau_n = 1 / f_n\n"
        f"tau_n = 1 / {round(f_n, precision)}\n"
        f"tau_n = {round(tau_n, precision)} s\n\n"

        f"**Answer:**\n"
        f"The undamped torsional natural frequency is {round(omega_n, precision)} rad/s or {round(f_n, precision)} Hz.\n"
        f"The natural period of vibration is {round(tau_n, precision)} seconds."
    )

    return question, solution
\end{lstlisting}


\vspace{6pt}

\subsection{Intermediate}

\vspace{6pt}

\begin{lstlisting}[style=pythonstyle]
def template_exponential_mttf_topology():
    configuration = random.choice(["series", "parallel"])

    if configuration == "series":
        lam1 = random.randint(*_T9_SERIES_RATE_CENTS) / 100
        lam2 = random.randint(*_T9_SERIES_RATE_CENTS) / 100
        lam3 = random.randint(*_T9_SERIES_RATE_CENTS) / 100
        lams = round(lam1 + lam2 + lam3, 2)          # exact 2-dp sum
        mttf = float((Decimal("1000") / Decimal(f"{lams:.2f}"))
                     .quantize(Decimal("0.1"), rounding=ROUND_HALF_UP))
        assert 333.0 <= mttf <= 1667.0, f"series MTTF out of bounds: {mttf}"

        question = (
            f"A booster station moves slurry through three pumps "
            f"connected in series along one line, so the station fails "
            f"as soon as ANY one pump fails. The pumps fail "
            f"independently, with exponentially distributed lifetimes "
            f"and constant failure rates of {lam1:.2f}, {lam2:.2f}, and "
            f"{lam3:.2f} failures per 1000 hours. Determine the mean "
            f"time to failure (MTTF -- the expected operating time until "
            f"the station first fails) of the station, in hours, to one "
            f"decimal place (round half up)."
        )
        solution = (
            f"**Given:**\n"
            f"Three pumps in series; independent exponential lifetimes "
            f"with failure rates lam1 = {lam1:.2f}, lam2 = {lam2:.2f}, "
            f"lam3 = {lam3:.2f} per 1000 hours.\n\n"
            f"**Step 1:** Characterize the station lifetime. The station "
            f"fails at the FIRST pump failure, i.e., at the minimum of "
            f"the three lifetimes; the minimum of independent "
            f"exponential lifetimes is exponential with the SUM of the "
            f"rates.\n\n"
            f"**Step 2:** Sum the failure rates.\n"
            f"lams = {lam1:.2f} + {lam2:.2f} + {lam3:.2f} = {lams:.2f} "
            f"failures per 1000 hours\n\n"
            f"**Step 3:** Invert the rate and convert to hours. A rate "
            f"of {lams:.2f} per 1000 hours means a mean lifetime of\n"
            f"MTTF = 1000 / {lams:.2f} = {mttf:.1f} hours (rounded half "
            f"up to one decimal)\n\n"
            f"**Answer:** The mean time to failure of the booster "
            f"station is {mttf:.1f} hours"
        )

    else:
        theta = random.randint(*_T9_THETA_RANGE)
        t_first = theta / 2                          # exact (x.0 or x.5)
        mttf = theta + t_first                       # 1.5*theta, exact
        assert 1500.0 <= mttf <= 7500.0, f"parallel MTTF out of bounds: {mttf}"

        question = (
            f"A cooling loop is served by two identical pumps in "
            f"parallel; the loop keeps operating as long as at least one "
            f"pump operates, and fails only when BOTH pumps have "
            f"failed. The pumps fail independently, with exponentially "
            f"distributed lifetimes averaging {theta} hours each. Using "
            f"the memoryless property, determine the mean time to "
            f"failure (MTTF -- the expected operating time until the "
            f"loop fails) of the cooling loop, in hours, to one decimal "
            f"place."
        )
        solution = (
            f"**Given:**\n"
            f"Two identical pumps in parallel; independent exponential "
            f"lifetimes with mean theta = {theta} hours (failure rate "
            f"1/{theta} per hour each).\n\n"
            f"**Step 1:** Split the loop lifetime into two phases: the "
            f"time until the FIRST pump failure, then the remaining "
            f"lifetime of the surviving pump.\n\n"
            f"**Step 2:** Expected time to the first failure. While both "
            f"pumps run, failures compete: the minimum of two "
            f"independent exponentials is exponential with doubled rate "
            f"2/{theta}, so\n"
            f"E[first failure] = theta / 2 = {theta} / 2 = {t_first:.1f} "
            f"hours\n\n"
            f"**Step 3:** Expected remaining lifetime of the survivor. "
            f"By the memoryless property, the surviving pump is "
            f"statistically good as new at the moment the other fails, "
            f"so its expected remaining lifetime is the full mean:\n"
            f"E[remaining] = theta = {theta} hours\n\n"
            f"**Step 4:** Add the phases.\n"
            f"MTTF = {t_first:.1f} + {theta} = {mttf:.1f} hours\n\n"
            f"**Answer:** The mean time to failure of the cooling loop "
            f"is {mttf:.1f} hours"
        )

    return question, solution
\end{lstlisting}


\vspace{6pt}

\subsection{Advanced}

\vspace{6pt}

\begin{lstlisting}[style=pythonstyle, literate={²}{{$^2$}}1 {³}{{$^3$}}1 {·}{{$\cdot$}}1 {₀}{{$_0$}}1 {₁}{{$_1$}}1 {₂}{{$_2$}}1 {≈}{{$\approx$}}1]
def template_vdw_solve_for_volume():
    # 1. Parameterize the inputs with a Validation Loop
    # We want to ensure we generate a problem with 3 roots (the "Two-Phase" scenario)
    # as this is the most pedagogically valuable case for cubic EOS problems.

    R = 0.08314  # L·bar/(mol·K)

    # Loop until we find a set of inputs that produces 3 distinct real roots
    while True:
        try:
            substance_name = random.choice(list(CRITICAL_PROPERTIES.keys()))
            properties = CRITICAL_PROPERTIES[substance_name]
            Tc = properties["Tc"]
            Pc = properties["Pc"]

            # Choose T well below critical point to ensure the "loop" is wide enough
            Tr = random.uniform(0.65, 0.85)
            T = round(Tr * Tc, 2)

            # Estimate VdW saturation pressure to hit the 3-root region
            # Approx VdW vapor pressure: log10(Pr) ~ -3(1/Tr - 1)
            approx_Pr_sat = 10**(-3.0 * (1.0/Tr - 1.0))

            # Pick P close to this saturation pressure
            Pr = approx_Pr_sat * random.uniform(0.95, 1.05)
            P = round(Pr * Pc, 3) # Use 3 decimals for precision

            # 2. Perform the core calculation
            a = (27 * (R**2) * (Tc**2)) / (64 * Pc)
            b = (R * Tc) / (8 * Pc)

            # Coefficients of the cubic polynomial: V³ + c₂V² + c₁V + c₀ = 0
            c2 = -(b + R * T / P)
            c1 = a / P
            c0 = -(a * b) / P
            coeffs = [1, c2, c1, c0]

            roots = np.roots(coeffs)

            # Filter for positive, real roots
            tolerance = 1e-9
            positive_real_roots = sorted([
                root.real for root in roots if abs(root.imag) < tolerance and root.real > 0
            ])

            # If we found 3 roots, we have a valid problem. Break the loop.
            if len(positive_real_roots) == 3:
                V_f = positive_real_roots[0]
                V_g = positive_real_roots[-1]
                break

        except Exception:
            continue

    # 3. Generate the question and solution strings
    question = (
        f"A vessel of {substance_name} is held at a temperature of {T} K and a pressure of {P} bar. "
        f"Using the van der Waals equation of state, determine the possible molar volumes (L/mol) "
        f"predicted for the liquid and vapor phases.\n\n"
        f"The critical properties for {substance_name} are:\n"
        f"Tc = {Tc} K\nPc = {Pc} bar"
    )

    solution = (
        f"**Step 1:** Calculate the van der Waals parameters 'a' and 'b'.\n"
        f"a = (27 * R^2 * Tc^2) / (64 * Pc) = {round(a, 4)} L^2·bar/mol^2\n"
        f"b = (R * Tc) / (8 * Pc) = {round(b, 5)} L/mol\n\n"

        f"**Step 2:** Formulate the cubic equation: V^3 + c2*V^2 + c1*V + c0 = 0.\n"
        f"The coefficients are derived from the VdW equation:\n"
        f"c2 = -(b + RT/P) = {round(c2, 4)} L/mol\n"
        f"c1 = a/P = {round(c1, 4)} L^2/mol^2\n"
        f"c0 = -(ab)/P = {round(c0, 6)} L^3/mol^3\n\n"

        f"**Step 3:** Solve the polynomial for its roots.\n"
        f"Using a numerical solver, we find three real, positive roots:\n"
        f"Roots: {', '.join([f'{r:.4f}' for r in positive_real_roots])} L/mol\n\n"
        f"\n\n"

        f"**Step 4:** Interpret the physical significance of the roots.\n"
        f"For a state in the subcritical region, the EOS predicts three roots:\n"
        f"- The **smallest root** corresponds to the molar volume of the **Liquid Phase**.\n"
        f"- The **largest root** corresponds to the molar volume of the **Vapor Phase**.\n"
        f"- The intermediate root is physically unstable and disregarded.\n\n"

        f"**Answer:**\n"
        f"1. Liquid-phase Volume (V_liq) ≈ **{round(V_f, 4)} L/mol**\n"
        f"2. Vapor-phase Volume (V_vap) ≈ **{round(V_g, 4)} L/mol**"
    )

    return question, solution
\end{lstlisting}
